\section{USAMO 1977}
        \begin{enumerate}
        	\item (\textit{USAMO 1977}) Determine all pairs of positive integers $(m,n)$ such that
$(1+x^n+x^{2n}+\cdots+x^{mn})$ is divisible by $(1+x+x^2+\cdots+x^{m})$.


 



\textbf{Solution 1}

Denote the first and larger polynomial to be $f(x)$ and the second one to be $g(x)$. In order for $f(x)$ to be divisible by $g(x)$ they must have the same roots. The roots of $g(x)$ are the (m+1)th roots of unity, except for 1. When plugging into $f(x)$, the root of unity is a root of $f(x)$ if and only if the terms $x^n, x^{2n}, x^{3n}, \cdots x^{mn}$ all represent a different (m+1)th root of unity not equal to 1. 

 Note that if $\\gcd(m+1,n)=1$, the numbers $n, 2n, 3n, \cdots, mn$ represent a complete set of residues minus 0 modulo $m+1$. However, if $gcd(m+1,n)=a$ not equal to 1, then $\frac{(m+1)(n)}{a}$ is congruent to $0 \pmod {m+1}$ and thus a complete set is not formed. Therefore, $f(x)$ divides $g(x)$ if and only if $\boxed{\\gcd(m+1,n)=1}.$ $\blacksquare$

 



\textbf{Solution 2}

We could instead consider $f(x)$ modulo $g(x)$. Notice that $x^{m+1} = 1 \pmod {g(x)}$, and thus we can reduce the exponents of $f(x)$ to their equivalent modulo $m+1$. We want the resulting $h(x)$ with degree less than $m+1$ to be equal to $g(x)$ (of degree $m$), which implies that the exponents of $f(x)$ must be all different modulo $m+1$. This can only occur if and only if $gcd(m+1, n) = 1$, and this is our answer, as shown in Solution 1.

 



\textbf{Solution 3}

Notice that $1+x^n+x^{2n}+...+x^{mn}=\frac{x^{mn+n}-1}{x^n-1}$ and $1+x+x^2+...+x^m=\frac{x^{m+1}-1}{x-1}$, so it remains to prove that $(x^n-1)(x^{m+1}-1) \mid (x^{mn+n}-1)(x-1)$. It is clear that $(x^n-1) \mid (x^{mn+n}-1)$ and $(x^{m+1}-1) \mid (x^{mn+n}-1)$. For any $k \in \mathbb{N}$, we can use the fact that $x^k-1=\prod_{d \mid k} \Phi _d(x)$, where $\Phi _d(x)$ is the dth cyclotomic polynomial. If $\gcd(m+1,n) = d > 1$, then $x^{m+1}-1$ and $x^n-1$ share a common cyclotomic polynomial; namely, $\Phi _d(x)$. But since all the factors of $x^{mn+n}-1$ are distinct, $(x^{mn+n}-1)(x-1)$ cannot be divisible by $(x^n-1)(x^{m+1}-1)$. We find that $\gcd(m+1,n) = 1$ must be the solution, since the only shared polynomial is $\Phi _1(x)=x-1$, and we are done.

 


	\item (\textit{USAMO 1977}) $ABC$ and $A'B'C'$ are two triangles in the same plane such that the lines $AA',BB',CC'$ are mutually parallel. Let $[ABC]$ denote the area of triangle $ABC$ with an appropriate $\pm$ sign, etc.; prove that

 \[3([ABC] + [A'B'C']) = [AB'C'] + [BC'A'] + [CA'B'] + [A'BC] + [B'CA] + [C'AB].\]


 


	\item (\textit{USAMO 1977}) If $a$ and $b$ are two of the roots of $x^4+x^3-1=0$, prove that $ab$ is a root of $x^6+x^4+x^3-x^2-1=0$.


 



\textbf{Solution}

Given the roots $a,b,c,d$ of the equation $x^{4}+x^{3}-1=0$.

 First, Vieta's relations give $a+b+c+d = -1 , ab+ac+ad+bc+bd+cd=0, abc+abd+acd+bcd=0, abcd = -1$.

 Then $cd=-\frac{1}{ab}$ and $c+d=-1-(a+b)$.

 The other coefficients give $ab+(a+b)(c+d)+cd = 0$ or $ab+(a+b)[-1-(a+b)]-\frac{1}{ab}=0$.

 Let $a+b=s$ and $ab=p$.

 Thus, $0=ab+ac+ad+bc+bd+cd=p+s(-1-s)-\frac{1}{p}$.   (1)

 Also, $0=abc+abd+acd+bcd=p(-1-s)-s/p$.

 Solving this equation for $s$, $s= \frac{-p^2}{p^2+1}$.

 Substituting into (1): $\frac{p^{6}+p^{4}+p^{3}-p^{2}-1}{p(p^2+1)^2}=0$.

 Conclusion: $p =ab$ is a root of $x^{6}+x^{4}+x^{3}-x^{2}-1=0$.

 


	\item (\textit{USAMO 1977}) Prove that if the opposite sides of a skew (non-planar) quadrilateral are congruent, then the line joining the midpoints of the two diagonals is perpendicular to these diagonals, and conversely, if the line joining the midpoints of the two diagonals of a skew quadrilateral is perpendicular to these diagonals, then the opposite sides of the quadrilateral are congruent.


 



\textbf{Solution}

We first prove that if the opposite sides are congruent, then the line joining the midpoints of the two diagonals is perpendicular to these diagonals.

 Let the vertices of the quadrilateral be A, B, C, D, in that order. Thus, the opposite sides congruent condition translates to (after squaring):

 
 \[AB^2 = CD^2, BC^2 = AD^2\]

 Let the position vectors of A, B, C, D be a, b, c, d with respect to an arbitrary origin $O$. Define the square of a vector $a^2 = a \cdot a$. Thus, $(a-b)^2 = (c-d)^2$ and so 
 \[a^2 - 2ab + b^2 = c^2 - 2cd + d^2.\] Similarly, 
 \[b^2 - 2bc + c^2 = a^2 - 2ad + d^2.\] Subtracting the second equation from the first eventually results in 
 \[a^2 - ab + cd + bc - ad - c^2 = 0,\] which factors as 
 \[(a - b + c - d)(a - c) = 0.\]

 Let M and N be the midpoints of $AC$ and $BD$, with position vectors $m = \frac{a+c}{2}$ and $n = \frac{b+d}{2}$, respectively. Then notice that the perpendicularity condition for vectors gives $MN$ perpendicular to $AC$. Similarly (by adding the equations), we find that $MN$ is perpendicular to $BD$, completing the first part of the proof.

 Proving the converse is straightforward also: indeed, $MN$ perpendicular to both diagonals gives

 
 \[a^2 - ab + cd + bc - ad - c^2 = 0\]
 \[b^2 - ab + cd - bc + ad - d^2 = 0\]

 Adding and subtracting the equations eventually simplifies to

 
 \[(a-b)^2 = (c-d)^2\]
 \[(a-d)^2 = (b-c)^2,\]

 or $AB^2 = CD^2$ and $AD^2 = BC^2$. This reduces to $AB = CD$ and $AD = BC$, as desired. $\blacksquare$

 



\textbf{Solution 2}

Let $A, B, C, D$ be position vectors rather than the specified points. We're given that $|A-B| = |C-D| \rightarrow |A-B|^2=|C-D|^2 \rightarrow (A-B) \cdot (A-B) = (C-D) \cdot (C-D) \rightarrow \\ A \cdot A - 2A \cdot B + B \cdot B - C \cdot C + 2C \cdot D - D \cdot D=0$. We're also given that $|A-D| = |C-B| \rightarrow |A-D|^2=|C-B|^2 \rightarrow (A-D) \cdot (A-D) = (C-B) \cdot (C-B) \rightarrow \\ A \cdot A - 2A \cdot D + D \cdot D - C \cdot C + 2C \cdot B - B \cdot B=0$. Adding the two equations gives $2A \cdot A - 2A \cdot B - 2A \cdot D - 2 C \cdot C + 2C \cdot D + 2C \cdot B = 0$. Dividing by 4 and rearranging gives $\frac{A+C-B-D}{2} \cdot (A-C) = 0$. On the left hand side of this equation, one of the two factors is the vector between the midpoints of the two diagonals, and the other factor is on of the two diagonals. Their dot product is $0$, so they are perpendicular. The same methods can be used to show that the vector between the two midpoints is perpendicular to the other diagonal. We can employ this process in reverse to get the converse (this isn't completely simple, but I'm lazy so it's left as an exercise i guess).

 


	\item (\textit{USAMO 1977}) If $a,b,c,d,e$ are positive numbers bounded by $p$ and $q$, i.e, if they lie in $[p,q], 0 < p$, prove that

 \[(a + b + c + d + e)\biggl(\frac {1}{a} + \frac {1}{b} + \frac {1}{c} + \frac {1}{d} + \frac {1}{e}\biggr) \le 25 + 6\left(\sqrt {\frac {p}{q}} - \sqrt {\frac {q}{p}}\right)^2\]
and determine when there is equality.


 



\textbf{Solution}

Fix four of the variables and allow the other to vary. Suppose, for example, we fix all but $x$. Then the expression on the LHS has the form $(r + x)(s + \frac{1}{x}) = (rs + 1) + sx + \frac{r}{x}$, where $r$ and $s$ are fixed. But this is convex. That is to say, as $x$ increases if first decreases, then increases. So its maximum must occur at $x = p$ or $x = q$. This is true for each variable.

 Suppose all five are $p$ or all five are $q$, then the LHS is 25, so the inequality is true and strict unless $p = q$. If four are $p$ and one is $q$, then the LHS is $17 + 4\left(\frac{p}{q} + \frac{q}{p}\right)$. Similarly if four are $q$ and one is $p$. If three are $p$ and two are $q$, then the LHS is $13 + 6\left(\frac{p}{q} + \frac{q}{p}\right)$. Similarly if three are $q$ and two are $p$.

 $\frac{p}{q} + \frac{q}{p} \geq 2$ with equality iff $p = q$, so if $p < q$, then three of one and two of the other gives a larger LHS than four of one and one of the other. Finally, we note that the RHS is in fact $13 + 6\left(\frac{p}{q} + \frac{q}{p}\right)$, so the inequality is true with equality iff either (1) $p = q$ or (2) three of $v, w, x, y, z$ are $p$ and two are $q$ or vice versa.

 


\end{enumerate}